\section{See the clipped partition}\label{transpose-x-tile-map}
\begin{figure}[htbp]
\centering
\begin{tikzpicture}[alt={A three-row five-column input matrix contains indices zero through fourteen. Tile A encloses columns zero through three in all three rows. Tile B encloses only column four, containing indices four, nine, and fourteen.},x=15mm,y=12mm]
\foreach \r in {0,1,2} {
 \foreach \c in {0,1,2,3,4} {
  \draw[gray] (\c,-\r) rectangle +(1,-1);
  \pgfmathtruncatemacro{\cell}{5*\r+\c}
  \node at (\c+0.5,-\r-0.5) {\cell};
 }
}
\draw[very thick] (0,0) rectangle (4,-3);
\draw[very thick,dashed] (4,0) rectangle (5,-3);
\node at (2,0.5) {Tile A: three by four};
\node[anchor=west,align=left] at (5.3,-1.5) {Tile B:\\three by one};
\end{tikzpicture}
\caption{The first tile is clipped in height; the second is clipped in both height and width. Labels and line styles distinguish the tiles without relying on color.}
\end{figure}
For a row-major traversal inside tile A, the source and destination word indices are
\[
\begin{aligned}
 \text{Source:}\quad&0,1,2,3,5,6,7,8,10,11,12,13,\\
 \text{Destination:}\quad&0,3,6,9,1,4,7,10,2,5,8,11.
\end{aligned}
\]
Tile B then fills output indices \(12,13,14\).
Together the two destination sets partition the fifteen output words. There is no
required sixteen-word temporary array for an edge tile, and no padding cell is
part of this traversal.
